\documentclass[11pt,a4paper,twoside,openright]{book}
\usepackage[english]{babel}
\newcommand{\HandbookPdfTitle}{SASD Math Toolkit - CSharp/.NET Handbook Version 1.0 RC1}
\newcommand{\HandbookFooter}{Handbook Version 1.0 RC1 - English}
\newcommand{\HandbookSubtitle}{Numerical Methods - C\# / .NET}
\newcommand{\HandbookVersionText}{User and Developer Handbook - Version 1.0 RC1}
\newcommand{\HandbookTagline}{Numerical methods for robust, inspectable and reusable C\# applications.}
\newcommand{\HandbookRepositoryRef}{Repository: \repo\quad |\quad C\#/.NET 10}

\usepackage{sasd-handbook}

\begin{document}
\frontmatter
\SASDTitlePage

\chapter*{Document status}
\addcontentsline{toc}{chapter}{Document status}
\begin{sasdnote}{Handbook Version 1.0 RC1}
This document is the English release-candidate handbook for SASD Math Toolkit package candidate \texttt{1.0.0-rc.1}. During the build, chapter content is generated from the maintained Markdown handbook under \texttt{docs/user-guide/en/}. Markdown remains the editorial source of truth while LaTeX provides typography, navigation and PDF output.
\end{sasdnote}

\begin{tabularx}{\textwidth}{@{}>{\bfseries}p{0.31\textwidth}X@{}}
Document & SASD Math Toolkit - C\#/.NET User and Developer Handbook\\
Handbook version & 1.0 RC1\\
Package candidate & \texttt{1.0.0-rc.1}\\
Repository & \repo\\
Target platform & .NET 10\\
Assembly & \texttt{Sasd.Math.Toolkit}\\
Root namespace & \texttt{Sasd.Numerics}\\
License & MIT\\
Author / project & Robin Goerlach; SASD-GmbH\\
\end{tabularx}

\chapter*{Preface}
\addcontentsline{toc}{chapter}{Preface}
SASD Math Toolkit carries the functional core of a classical numerical-methods toolbox into a modern, independently developed C\#/.NET library and then extends that foundation for current 2026 workloads. The historical manual is used as orientation for problem classes and handbook structure, not as a source-text or source-code template.

Chapters 1--11 describe the classical numerical foundation. Chapters 12 and 13 document the modern dense and sparse linear-algebra layers added during the 2026 modernization: Householder QR, Cholesky, SVD, conditioning diagnostics, canonical CSR, preconditioning, CG/PCG, restarted GMRES and BiCGSTAB.

\begin{sasdtip}{Guiding principle}
A number alone is not yet a trustworthy numerical result. For important computations, method, tolerance, termination status and a problem-specific quality check belong together.
\end{sasdtip}

\tableofcontents
\clearpage
\mainmatter

\input{build/en/getting-started.tex}
\input{build/en/root-finding.tex}
\input{build/en/interpolation.tex}
\input{build/en/differentiation.tex}
\input{build/en/integration.tex}
\input{build/en/matrices-linear-systems.tex}
\input{build/en/eigenvalues-eigenvectors.tex}
\input{build/en/ordinary-differential-equations.tex}
\input{build/en/least-squares.tex}
\input{build/en/fft-convolution-correlation.tex}
\input{build/en/diagnostics-common-failure-modes.tex}
\input{build/en/modern-dense-linear-algebra.tex}
\input{build/en/sparse-linear-algebra.tex}

\appendix
\chapter{Historical context and current V1 coverage}
The C\# edition deliberately follows the major numerical problem classes of the historical Numerical Methods Toolbox at a high level. It is not a source-code port. Public APIs, data types, error handling, examples and diagnostics belong to SASD Math Toolkit. The 2026 modernization goes beyond that historical catalog where current SASD applications need stronger numerical foundations.

\begin{longtable}{@{}p{0.28\textwidth}p{0.62\textwidth}@{}}
\toprule
\textbf{Topic} & \textbf{C\#/.NET edition}\\
\midrule
Root finding & Bisection, Newton-Raphson, secant, Newton-Horner, Muller, Laguerre and polynomial deflation\\
Interpolation & Lagrange, divided differences, natural and clamped cubic splines\\
Differentiation / integration & Function, tabular and spline differentiation; Newton-Cotes, adaptive rules, Romberg and Gauss-Legendre\\
Dense matrix routines & \texttt{DenseMatrix}, Gaussian elimination, pivoted LU, Householder QR, Cholesky, SVD, pseudoinverse, norms, conditioning, inverse, residuals and Gauss-Seidel\\
Sparse matrix routines & Canonical CSR, sparse matvec/norms, Jacobi preconditioning, CG/PCG, restarted GMRES, BiCGSTAB and independent true-residual diagnostics\\
Eigenvalue problems & Power/inverse-power methods, Wielandt deflation and cyclic Jacobi\\
Differential equations & RK4 family, RKF45, Adams predictor-corrector, linear and nonlinear shooting\\
Least squares & Polynomial/general-basis fits and named models; Householder QR for full-column-rank problems with SVD fallback for rank-deficient or underdetermined minimum-norm problems\\
FFT / applications & Complex and real FFT, compact real spectrum, convolution and cross-correlation\\
Demonstrations & A platform-neutral deterministic HTML/SVG report plus a modern sparse release smoke solve\\
\bottomrule
\end{longtable}

\chapter{Building and maintaining the LaTeX edition}
The LaTeX edition is intentionally generated from the existing Markdown chapters. Editorial changes should therefore be made first in \texttt{docs/user-guide/de/} or \texttt{docs/user-guide/en/}. The \texttt{LaTeX/} directory owns layout, build logic and language-specific front/back matter.

\begin{lstlisting}
cd docs/user-guide/LaTeX
make de      # German PDF
make en      # English PDF
make all     # both editions
\end{lstlisting}

\backmatter
\chapter*{License and provenance}
\addcontentsline{toc}{chapter}{License and provenance}
SASD Math Toolkit is licensed under MIT. Historical Borland material and historical product names remain the property of their respective rights holders. This handbook documents the independently developed SASD C\#/.NET implementation and does not reproduce the historical manual text.

\end{document}
